\clearpage
\VerdierPhysicalPageAnchor{0146}%
\begingroup
\small
\setlength{\parskip}{0pt}%
\setlength{\emergencystretch}{1.5em}%
\setlength{\abovedisplayskip}{3pt plus 1pt minus 1pt}%
\setlength{\belowdisplayskip}{3pt plus 1pt minus 1pt}%

\begingroup\itshape
\VerdierSemanticAnchor{ch2-2-3-3-f}
\noindent f) An object $Z$ of $\mathcal D$ belongs to
$Q^{-1}(\operatorname{Lib}_R(Q))$ if and only if there exist an object $Y$
free on the right and a morphism $s\colon Z\longrightarrow Y$, where $s$
belongs to $\widehat{S_{\mathcal D}(\mathcal B)}$, the saturation of the
system $S_{\mathcal D}(\mathcal B)$ (2.2.4, (c)).

\VerdierSemanticAnchor{ch2-2-3-3-g}
\par\noindent g) The inclusion functor
$L_R(Q)\subset Q^{-1}(\operatorname{Lib}_R(Q))$ admits a left adjoint
\[
  \mathcal R\colon Q^{-1}(\operatorname{Lib}_R(Q))\longrightarrow L_R(Q).
\]

\VerdierSemanticAnchor{ch2-2-3-3-adjunction-diagram}
\begin{center}
\begin{tikzcd}[column sep=3.2em,row sep=3.1em]
L_R(Q)
  \arrow[r,shift left=.7ex]
  \arrow[dr,"{Q|_{L_R(Q)}}"']
& Q^{-1}(\operatorname{Lib}_R(Q))
  \arrow[l,dashed,bend right=28,"{\mathcal R}"']
  \arrow[r,hook]
  \arrow[d,"{Q_1}"]
& \mathcal D \arrow[d,"Q"] \\
& \operatorname{Lib}_R(Q)
  \arrow[ul,dashed,bend left=10,"{\Psi}"']
  \arrow[r,hook]
& \mathcal D/\mathcal B
\end{tikzcd}
\end{center}
\endgroup

\VerdierSemanticAnchor{ch2-2-3-3-adjunction-explanation}
Let $\Psi\colon\operatorname{Lib}_R(Q)\longrightarrow L_R(Q)$ be a
quasi-inverse to $Q|_{L_R(Q)}$ (cf. c) and e)). We may take
$\mathcal R=\Psi Q_1$, where
$Q_1\colon Q^{-1}(\operatorname{Lib}_R(Q))\longrightarrow
\operatorname{Lib}_R(Q)$ denotes the functor induced by $Q$. The functor
$\Psi$ is right adjoint to the functor $Q_1$.

\VerdierSemanticAnchor{ch2-2-3-4}
\medskip
\noindent\textbf{2.3.4. Proof of Proposition 2.3.3:}\quad
\VerdierSemanticAnchor{ch2-2-3-4-a-proof}
Let us prove a). First, if $Y$ is $Q$-free on the right, the functor
$X\longmapsto\operatorname{Hom}_{\mathcal D}(X,Y)$ is isomorphic to the
functor
$X\longmapsto\operatorname{Hom}_{\mathcal D/\mathcal B}(Q(X),Z)$, for a
suitable object $Z$ of $\mathcal D/\mathcal B$. Consequently, the functor
$X\longmapsto\operatorname{Hom}_{\mathcal D}(X,Y)$ sends the morphisms in
$S_{\mathcal D}(\mathcal B)$ to isomorphisms. This proves i)
$\Rightarrow$ ii). The implication ii) $\Rightarrow$ iii) follows from the
fact that the morphism $B\longrightarrow0$ belongs to
$S_{\mathcal D}(\mathcal B)$. The implication iii) $\Rightarrow$ iv)
follows from the definition of $S_{\mathcal D}(\mathcal B)$ and (1.2.6).
The implication iv) $\Rightarrow$ v) follows from the isomorphism (2.2.4):

\VerdierSemanticAnchor{ch2-2-3-4-colimit-isomorphism}
\[
  \operatorname{Hom}_{\mathcal D/\mathcal B}(Q(X),Q(Y))
  \xrightarrow{\sim}
  \varinjlim_{Y\backslash S_{\mathcal D}(\mathcal B)}
  \operatorname{Hom}_{\mathcal D}(X,\mathord\cdot)\, .
\]

\VerdierSemanticAnchor{ch2-2-3-4-a-proof-complete}
Finally, the implication v) $\Rightarrow$ i) is evident: $Y$ is free on
the right over $Q(Y)$.

\VerdierSemanticAnchor{ch2-2-3-4-b-c-e-proof}
Assertion b) follows from a), iii). Assertion c) follows from a), v). It
also follows from a), v) that the objects of $\mathcal D/\mathcal B$ which
are liberable on the right are precisely the objects of
$\mathcal D/\mathcal B$ isomorphic in $\mathcal D/\mathcal B$ to images
under $Q$ of objects that are $Q$-free on the right. The freedom arrows are
isomorphisms. This proves e). It is then clear that
$\operatorname{Lib}_R(Q)$ is a strictly full triangulated subcategory.

\VerdierSemanticAnchor{ch2-2-3-4-d-proof}
Let $Z=Z_1\oplus Z_2$ be an object of $\operatorname{Lib}_R(Q)$. There is
therefore an object $Y$ that is $Q$-free on the right and an isomorphism
$Q(Y)\xrightarrow{\sim}Z_1\oplus Z_2$. The functor
$X\longmapsto\operatorname{Hom}_{\mathcal D}(X,Y)$ therefore admits a
projector, and hence the decomposition of $Z$ as a direct sum defines a
projector on $Y$. The factors of $Z$ will belong to
$\operatorname{Lib}_R(Q)$ if and only if this projector is decomposable,
which is the case when countable sums or products are representable in
$\mathcal D$ (1.2.9).
\endgroup
